\documentclass[a4paper,10pt]{article} %Adjust Font Size Below NOT here 

% Packages
\usepackage[T1]{fontenc}
\usepackage{XCharter}
\usepackage[margin=0.4in]{geometry}
\usepackage{enumitem}
\usepackage[hidelinks]{hyperref}
\usepackage{titlesec}
\usepackage{setspace} % for line spacing
\usepackage{fontawesome}
\usepackage[fontsize=11  pt]{scrextend}

% Formatting
\titleformat{\section}{\large\bfseries}{}{0em}{}[\titlerule]
\setlist[itemize]{noitemsep, topsep=0pt}
\renewcommand{\baselinestretch}{1.1} % increase line spacing
\tolerance=1 \emergencystretch=\maxdimen \hyphenpenalty=10000 \hbadness=10000
\begin{document}
\pagestyle{empty}
\begin{center}
    {\huge\textbf{Amit Chandran}} \\
    \href{mailto:amitchandran06@gmail.com}{amitchandran06@gmail.com} \textbar{} \href{https://www.linkedin.com/in/amit-chandran}{LinkedIn \faLinkedinSquare} \textbar{} \href{https://github.com/amitchandran06}{GitHub \faGithub} \textbar{} \href{https://amitchandran06.github.io/Portfolio/}{Portfolio \faGlobe} \\
    \vspace{0.25cm}
    
        {\textrm{Multidisciplinary systems engineer with experience developing C++ firmware, low-latency telemetry pipelines, and fully integrated electro-mechanical systems. Proven ability to deliver high-precision software and hardware solutions, including leveraging AI developer tools safely and effectively.}}
    
\end{center}
\vspace{-0.75cm}
\section*{Education}
\noindent\textbf{UCL - MEng Mechanical Engineering} \hfill \textbf{Expected June 2028}
\begin{itemize}
    \item Relevant Modules: Machine Learning, Dynamics, Control Systems
    \item Current Average: 78\% - First Class
\end{itemize}

\vspace{-0.5cm}
\section*{Experience}
\noindent\textbf{Systems Engineering Intern – Leonardo UK} \hfill \textbf{Summer 2025 and 2026}
\begin{itemize}
    \item Coordinated 10 interns to integrate antenna, firmware, hardware, and software into a continuous RF-to-web telemetry pipeline
    \item Developed a compression algorithm reducing FFT data bandwidth requirements by over \textbf{2500×}, minimising network utilisation whilst preserving critical information
    \item Automated manual RF signal analysis in MATLAB and reduced complexity from $\mathcal{O}(N^2)$ to $\mathcal{O}(N)$, cutting simulation time by \textbf{99.3\%}
    \item Launched a student-led cross-functional collaboration initiative, now incorporated into the internship program for future cohorts
\end{itemize}
\vspace{-0.75em}
\section*{Projects / Extracurricular}
\noindent\textbf{Quadruped QDD Actuator Driver} \hfill \textbf{June 2026 – Present}
\begin{itemize}
    \item Designed a bespoke multi-layer PCB with an onboard STM32 and 3-phase inverter to drive robotic actuators, achieving a cost per unit around \textbf{70\%} lower than commercial alternatives
    \item Developed custom C++ control algorithms for the STM32 to execute closed-loop \textbf{20~kHz} Field-Oriented Control (FOC), providing a modular framework for precise torque and position control
    \item Integrated the CAN protocol into the firmware, allowing for rapid telemetry exchange across 12 drivers
    \item Architected a multi-tiered safety system combining hardware-interrupt fault detection and SEGGER RTT logging to safely isolate overcurrent, undervoltage, and CAN  faults
     
\end{itemize}
    
    
\vspace{0.75em}
\noindent\textbf{Technical Lead and Drone Controls Engineer – UCL Rover Team} \hfill \textbf{January 2025 – Present}
\begin{itemize}
    \item Programmed an edge-computing vision pipeline in Python and OpenCV on a Raspberry Pi to track ArUco markers in real-time, with an average pose error of \textbf{$<$1cm}
    \item Established a robust UDP telemetry link to stream positional data seamlessly from the edge device to a high-level MATLAB/Simulink control model
    \item Built a UART-based MAVLink telemetry architecture connecting Simulink models directly to an STM32H743 flight controller for real-time flight control
\end{itemize}
\vspace{0.75em}
\noindent\textbf{Smart Pet Feeder – BLE IoT System} \hfill \textbf{October 2025 – March 2026}
\begin{itemize}
    \item Developed a smart pet feeder using ESP32-S3 C++ firmware, combining a stepper motor and load cell to dispense food portions with \textbf{±1\% mass accuracy} in a closed-loop system
    \item Implemented an event-driven BLE GATT server (Nordic UART Service) for mobile time-synchronization, eliminating external RTC costs; open-sourced firmware and DfMA CAD models
    \item Interfaced the feeder with an Android app via BLE, enabling remote feeding and real-time weight monitoring
\end{itemize}


\section*{Skills}
\begin{itemize}
    \item \textbf{Programming Languages} – C, C++, Python, MATLAB
    \item \textbf{Tools / Software} – EasyEDA (PCB), Git, Simulink, Fusion 360, Azure DevOps
    \item \textbf{Manufacturing / Lab} – Oscilloscopes, Soldering / SMD Rework, 3D Printing, Laser Cutting
\end{itemize}

\end{document}
